% ============================================================
%  第二卷 第十三章 引力波（§107–§110 + 习题）
%  底本：高教社中文版第8版；OCR 错误已修正，脚注文本待脚注代理补录
%  主要修正：
%  - OCR 将达朗贝尔算符 □ 排为 ☐/▽/⊔ 等，统一为 \square；
%    克里斯托费尔符号 Γ 排为 𝒯，恢复为 \Gamma；𝔤^{ik}（γ-密度）用 \mathfrak{g}；
%  - §107：(107.10) 之后的文字"分量ψ1,ψ2,ψ3,ψ0也化为零"按 (107.10) 补全指标
%    为 ψ_1^1, ψ_2^1, ψ_3^1, ψ_0^1；
%  - §108：(108.3) 后两项指标序排乱，按 (108.5) 同一组合复原；
%  - §109："弱平面引力波的推 r²(1)"为"推广①"之误；"-g~η^{4①}"圈码移出指数；
%    习题"f(t-x,y,z)=0 应该是调和函数"衍 "=0" 删；"线性关☉"→"线性关系"；
%  - §110：(110.14) 偏振平均公式中乱排的 "=" 删（该式为单一等式）；
%    习题1 "在圆周运动时 r=const, ψ̇=…" 烂字 κ̄ 复原为 ψ̇；
%    习题2 作者后 OCR 圈码 ① 保留；
%  - §110 习题3：能量式的 D^{(V)}（五点导数）与力、(3) 式的导数阶数
%    按 OCR 与习题4 圆轨道极限核对复原（f 与 (3) 中导数阶数存疑，见交付报告）；
%  - 习题内公式保留原书局部编号（\tag 实现），正文公式自动编号；
%  - 脚注圈码①照 OCR 保留（§109 一处、§110 三处）。
%  - 本章无插图。
%  - 2026-10-10 脚注已转录：5 处圈码标记全部替换为 \footnote{...}（§109 两处、
%    §110 三处），文本转录自原书扫描页 printed p.390/391/392/394/396。
% ============================================================
\chapter{引力波}

\section{弱引力波}

就像电动力学中那样，相对论引力理论中相互作用的传播速度的有限性使得与物体没有联系的自由引力场——引力波的存在成为可能.

现在我们来研究真空中的弱自由引力场，如同§105，引入描述伽利略度规的微弱扰动的张量 $h_{ik}$：
\begin{equation}
  g_{ik}=g_{ik}^{(0)}+h_{ik}.
\end{equation}

于是，准确到 $h_{ik}$ 的一阶量，逆变度规张量是：
\begin{equation}
  g^{ik}=g^{ik(0)}-h^{ik},
\end{equation}
而张量 $g_{ik}$ 的行列式：
\begin{equation}
  g=g^{(0)}(1+h),
\end{equation}
其中 $h\equiv h_i^i$；所有升高和降低张量指标的运算都按照未经扰动的度规 $g_{ik}^{(0)}$ 进行.

就像§105中已经指出的那样，$h_{ik}$ 为小量的条件使得作形式为 $x'^i=x^i+\xi^i$（$\xi^i$ 为小量）的参考系的任意变换成为可能；在这种条件下：
\begin{equation}
  h'_{ik}=h_{ik}
  -\frac{\partial\xi_i}{\partial x^{k}}
  -\frac{\partial\xi_k}{\partial x^{i}}.
\end{equation}

利用张量 $h_{ik}$ 的这种规范任意性，我们对它加上补充条件
\begin{equation}
  \frac{\partial\psi_i^k}{\partial x^{k}}=0,
  \qquad
  \psi_i^k=h_i^k-\frac{1}{2}\delta_i^kh,
\end{equation}
在这之后里奇张量具有简单的形式 (105.11)：
\begin{equation}
  R_{ik}=\frac{1}{2}\square h_{ik},
\end{equation}
其中 $\square$ 表示达朗贝尔算符：
\begin{equation*}
  \square=-g^{lm(0)}
  \frac{\partial^{2}}{\partial x^{l}\partial x^{m}}
  =\Delta-\frac{1}{c^{2}}\frac{\partial^{2}}{\partial t^{2}}.
\end{equation*}

条件 (107.5) 仍旧不能唯一地确定参考系的选取：如果某些 $h_{ik}$ 满足这些条件，那么 (107.4) 式的 $h'_{ik}$ 也将满足它们，只要 $\xi^i$ 是下面方程的解：
\begin{equation}
  \square\xi^i=0.
\end{equation}

令表达式 (107.6) 等于零，这样我们求得如下形式的真空中的引力场方程：
\begin{equation}
  \square h_i^k=0.
\end{equation}

这就是普通的波动方程.因此，引力场也同电磁场一样，以光速在真空中传播.

我们来研究一个平面引力波.在这样的波内，场仅沿着空间的一个方向变化；我们选择坐标轴 $x^{1}=x$ 作为这个方向.方程 (107.8) 这时变为
\begin{equation}
  \left(\frac{\partial^{2}}{\partial x^{2}}
  -\frac{1}{c^{2}}\frac{\partial^{2}}{\partial t^{2}}\right)h_i^k=0,
\end{equation}
它的解是 $t\pm x/c$ 的任意函数（见§47）.

假设波向着 $x$ 轴的正方向传播.由此，所有的 $h_i^k$ 都是 $t-x/c$ 的函数.辅助条件 (107.5) 在这种情形下给出 $\dot{\psi}_i^1-\dot{\psi}_i^0=0$，此处符号上的一点表示对 $t$ 微分.这个等式可以简单地通过去掉微分符号就能积出，积分常数可以设为零，因为我们所感兴趣的只是场的可变部分（正如电磁波的情形一样）.因此，$\psi_i^k$ 的各分量之间有关系式
\begin{equation}
  \psi_1^1=\psi_1^0,\quad
  \psi_2^1=\psi_2^0,\quad
  \psi_3^1=\psi_3^0,\quad
  \psi_0^1=\psi_0^0.
\end{equation}

正如上文已经指出的，条件 (107.5) 也还不能唯一地决定参考系；我们可以给坐标施加 $x'^i=x^i+\xi^i(t-x/c)$ 形式的变换；这些变换可以用来使四个量 $\psi_1^0,\psi_2^0,\psi_3^0,\psi_2^2+\psi_3^3$ 化为零；从等式 (107.10) 可以推断，这时分量 $\psi_1^1,\psi_2^1,\psi_3^1,\psi_0^1$ 也化为零.至于余下的量 $\psi_2^3,\psi_2^2-\psi_3^3$，无论怎样选择参考系也不能化为零，因为从 (107.4) 可以看出，在带 $\xi_i=\xi_i(t-x/c)$ 的变换下，这些分量一般不改变.我们注意到 $\psi\equiv\psi_i^i$ 这时也为零，所以 $\psi_i^k=h_i^k$.

因此，平面引力波由 $h_{23},h_{22}=-h_{33}$ 两个量所决定.换句话说，引力波是横波，波的偏振为 $yz$ 平面内的二阶对称张量所决定，这个张量的对角线的分量之和 $h_{22}+h_{33}$ 为零.选取 $h_{23}$ 和 $(h_{22}-h_{33})/2$ 这两个量的其中之一不为零的两种情况，我们可以得到两个独立的偏振.这样的两个偏振之间的区别是在平面 $yz$ 内旋转 $\pi/4$ 的角度.

我们来计算平面引力波里的能动赝张量.分量 $t^{ik}$ 是二阶小量；我们必须在忽略更高阶项的条件下算出它们.由于当 $h=0$ 时行列式 $g$ 与 $g^{(0)}=-1$ 的区别仅仅是二阶量，所以在通式 (96.9) 中可以规定 $\mathfrak{g}^{ik}{}_{,l}\approx g^{ik}{}_{,l}\approx-h^{ik}{}_{,l}$.对于平面波而言，$t^{ik}$ 中的所有的非零项都包括在下面的项中
\begin{equation*}
  \frac{1}{2}g^{il}g^{km}g_{np}g_{qr}g^{nr}{}_{,l}g^{pq}{}_{,m}
  =\frac{1}{2}h_q^{n,i}h_n^{q,k},
\end{equation*}
该项包含在 (96.9) 式的花括号中（只要选择伽利略参考系的一条轴作为波的传播方向就能容易地证实这一点）.这样一来
\begin{equation}
  t^{ik}=\frac{c^{4}}{32\pi k}h_q^{n,i}h_n^{q,k}.
\end{equation}

波里的能流由物理量 $-cgt^{0\alpha}\approx ct^{0\alpha}$ 确定.在沿轴 $x^{1}$ 传播的平面波中，非零的量 $h_{23}$ 和 $h_{22}=-h_{33}$ 只依赖于 $t-x/c$，这个能流的方向沿着同一个轴 $x^{1}$ 并且等于
\begin{equation}
  ct^{01}=\frac{c^{3}}{16\pi k}\left[\dot{h}_{23}^{2}
  +\frac{1}{4}(\dot{h}_{22}-\dot{h}_{33})^{2}\right].
\end{equation}

任意引力波场的初始条件应该由坐标的四个任意函数给定：由于波的横波特性，总共只有两个独立的分量 $h_{\alpha\beta}$，除此之外还应该给定它们对时间的一阶导数.尽管我们在这里进行的计算是以弱引力场的性质为出发点，但很明显，它的结果——数目 4——不可能依赖于这些前提条件，并适用于任何自由的，即与引力质量无关的引力场.

\begin{xiti}

\noindent{\heiti 习题}\quad 确定弱平面引力波里的曲率张量.

解：按照公式 (105.8) 计算 $R_{iklm}$，得到下列不为零的分量：
\begin{equation*}
  -R_{0202}=R_{0303}=-R_{1212}=R_{0212}=R_{0331}=R_{3131}=\sigma,
\end{equation*}

\begin{equation*}
  R_{0203}=-R_{1231}=-R_{0312}=R_{0231}=\mu,
\end{equation*}
其中有下列符号表示：
\begin{equation*}
  \sigma=-\frac{1}{2}\ddot{h}_{33}=\frac{1}{2}\ddot{h}_{22},
  \qquad
  \mu=-\frac{1}{2}\ddot{h}_{23}.
\end{equation*}

利用 (92.15) 式引入的三维张量 $A_{\alpha\beta}$ 和 $B_{\alpha\beta}$，我们有
\begin{equation*}
  A_{\alpha\beta}=
  \begin{pmatrix}
    0 & 0 & 0\\
    0 & -\sigma & \mu\\
    0 & \mu & \sigma
  \end{pmatrix},
  \qquad
  B_{\alpha\beta}=
  \begin{pmatrix}
    0 & 0 & 0\\
    0 & \mu & \sigma\\
    0 & \sigma & -\mu
  \end{pmatrix}.
\end{equation*}

适当地旋转坐标轴 $x^{2},x^{3}$ 能够将 $\sigma$ 或 $\mu$ 的其中之一变为零（在四维空间中的给定的一点）；量 $\sigma$ 变为零时，我们就把曲率张量化为简并的彼得罗夫Ⅱ型（N 型）.

\end{xiti}

\section{弯曲时空内的引力波}

就像我们在平直时空的“背景”下研究引力波的传播那样，可以考察相对于任意（非伽利略的）“未扰动”度规 $g_{ik}^{(0)}$ 的微小扰动的传播.顾及到某些其他可能的应用，在这里我们将必要的公式写成最一般的形式.

再次将 $g_{ik}$ 写成 (107.1) 的形式，我们求得通过修正值 $h_{ik}$ 表达的克里斯托夫符号的一阶修正：
\begin{equation}
  \Gamma_{kl}^{i(1)}
  =\frac{1}{2}\left(h_{k;l}^{i}+h_{l;k}^{i}-h_{kl}{}^{;i}\right),
\end{equation}
这一点可以通过直接的计算得到确认（这里和下面所有的张量运算——升降指标，协变微分——都借助于非伽利略度规 $g_{ik}^{(0)}$ 进行）.对于曲率张量的修正值我们得到：
\begin{equation}
\begin{aligned}
  R_{klm}^{i(1)}
  =\frac{1}{2}\bigl(&h_{k;m;l}^{i}+h_{m;k;l}^{i}
  -h_{km}{}^{;i}{}_{;l}\\
  &-h_{k;l;m}^{i}-h_{l;k;m}^{i}+h_{kl}{}^{i}{}_{;m}\bigr).
\end{aligned}
\end{equation}

由此可得里奇张量的修正值：
\begin{equation}
  R_{ik}^{(1)}=R^{l(1)}{}_{ilk}
  =\frac{1}{2}\left(h_{i;k;l}^{l}+h_{k;i;l}^{l}
  -h_{ik}{}^{;l}{}_{;l}-h_{;i;k}\right).
\end{equation}

里奇张量的混合分量的修正值可以从下面的关系式得到：
\begin{equation*}
  R_i^{k(0)}+R_i^{k(1)}
  =(R_{il}^{(0)}+R_{il}^{(1)})(g^{kl(0)}-h^{kl}),
\end{equation*}
由此
\begin{equation}
  R_i^{k(1)}=g^{kl(0)}R_{il}^{(1)}-h^{kl}R_{il}^{(0)}.
\end{equation}

真空中的精确度规应该满足精确的爱因斯坦方程 $R_{ik}=0$.由于未扰动度规 $g_{ik}^{(0)}$ 满足方程 $R_{ik}^{(0)}=0$，那么对于扰动得到方程 $R_{ik}^{(1)}=0$，即
\begin{equation}
  h_{i;k;l}^{l}+h_{k;i;l}^{l}
  -h_{ik}{}^{;l}{}_{;l}-h_{;i;k}=0.
\end{equation}

在任意引力波的一般情况下，将这个方程简化到类似于 (107.8) 的形式是不可能的.然而在高频波这一重要情况下可以做到这一点：波长 $\lambda$ 和振动周期 $\lambda/c$ 与表征“背景场”变化的特征距离 $L$ 和特征时间 $L/c$ 相比是小量.和未扰动度规 $g_{ik}^{(0)}$ 的导数相比，分量 $h_{ik}$ 的每次微分都会提高一个因子 $L/\lambda$.如果将精度限定在两个最高阶的项（$(L/\lambda)^{2}$ 和 $(L/\lambda)$），那么在 (108.5) 中我们可以交换微分的顺序；实际上，差值
\begin{equation*}
  h_{i;k;l}^{l}-h_{i;l;k}^{l}
  \approx h_m^lR^{m(0)}{}_{ikl}
  -h_i^mR^{l(0)}{}_{mkl}
\end{equation*}
具有 $(L/\lambda)^{0}$ 阶，而表达式 $h_{i;k;l}^{l}$ 和 $h_{i;l;k}^{l}$ 中的每一个都包含两个更高阶的项.现在给 $h_{ik}$ 规定一个补充条件
\begin{equation}
  \psi_{i;k}^{k}=0
\end{equation}
（类似于 (107.5)），我们得到方程
\begin{equation}
  h_{ik}{}^{;l}{}_{;l}=0,
\end{equation}
它是方程 (107.8) 的推广.

根据在§107中指出的原因，条件 (108.6) 并未唯一确定坐标的选取.对后者仍可以施加变换 $x'^i=x^i+\xi^i$，其中小量 $\xi^i$ 满足方程 $\xi^{i;k}{}_{;k}=0$.这些变换可以特别用来给 $h_{ik}$ 规定条件 $h\equiv h_i^i=0$.那么 $\psi_i^k=h_i^k$，于是 $h_i^k$ 符合条件
\begin{equation}
  h_{i;k}^{k}=0,\qquad h=0.
\end{equation}

在这样的规定之后，容许的变换就归结为条件 $\xi^i{}_{;i}=0$.

一般说来，赝张量 $t^{ik}$ 除了包含未扰动部分 $t^{ik(0)}$，也包含 $h_{ik}$ 的各阶项.如果我们考察在四维空间的一些区域上求平均后的量 $t^{ik}$，并且这些区域的尺寸与 $\lambda$ 相比很大，而与 $L$ 相比很小，那么我们会得到类似于 (107.11) 的表达式.这样的平均（下面用尖括号表示 $\langle\cdots\rangle$）不会影响 $g_{ik}^{(0)}$，却会使得关于快速振荡量 $h_{ik}$ 的所有线性项变为零.二次项中我们只保留那些关于 $1/\lambda$ 的最高（二）阶项；这就是关于导数 $h_{ik,l}\equiv\partial h_{ik}/\partial x^{l}$ 的二次项.

在这样的精度下，$t^{ik}$ 中所有的表现为四维散度的项可以被忽略.实际上，对这样的表达式沿四维空间的区域（求平均的区域）的积分按照高斯定理进行变换，结果会导致它们关于 $1/\lambda$ 的量级减少 1.除此之外，在分部积分之后按照 (108.7) 和 (108.8) 式化为零的那些项也会消失.于是，进行分部积分并且忽略对四维散度的积分，我们得到：
\begin{equation*}
  \langle h^{ln}{}_{,p}h_{l,n}^{p}\rangle
  =-\langle h^{ln}h_{l,p,n}^{p}\rangle=0,
\end{equation*}

\begin{equation*}
  \langle h^{il}{}_{,n}h_{l}^{k,n}\rangle
  =-\langle h^{il}h_{l,n}^{k,n}\rangle=0.
\end{equation*}

结果从所有的二阶项中仅剩下
\begin{equation}
  \langle t^{ik(2)}\rangle
  =\frac{c^{4}}{32\pi k}\langle h_q^{n,i}h_n^{q,k}\rangle.
\end{equation}

我们指出，在这种情况下，以同样的精度，$\langle t_i^{i(2)}\rangle=0$.

引力波具有一定的能量，它本身成为某个附加引力场的源.这个场同产生它的能量一起是关于 $h_{ik}$ 的二阶效应.但在高频引力波的情况下这个效应有实质性的加强：事实上，赝张量 $t^{ik}$ 是 $h_{ik}$ 的导数的二次式，这会将一个大的因子 $\lambda^{-2}$ 带入它的量级中.在这种情况下可以说，引力波本身产生了背景场，它们就在这个背景场上传播.按照上面的描述，在四维空间内尺度远大于 $\lambda$ 的区域求平均后，可以方便地对这个场进行研究.这样的平均运算抹平了短波的“涟漪”，产生了缓慢变化的背景度规 (R. A. Isaacson, 1968).

为了推导确定这个度规的方程，在张量 $R_{ik}$ 的展开式中应该不仅要考虑线性项，还要考虑关于 $h_{ik}$ 的二次项：$R_{ik}=R_{ik}^{(0)}+R_{ik}^{(1)}+R_{ik}^{(2)}$.正如已经指出的那样，求平均运算不会影响到零阶项.这样一来，平均后的场方程 $\langle R_{ik}\rangle=0$ 具有如下形式：
\begin{equation}
  R_{ik}^{(0)}=-\langle R_{ik}^{(2)}\rangle,
\end{equation}
并且在 $R_{ik}^{(2)}$ 中只应保留关于 $1/\lambda$ 的二次项.它们可以容易地从恒等式 (96.7) 得到.这个恒等式的右边具有四维散度的形式，从这个恒等式右边产生的关于 $h_{ik}$ 的二次项在求平均的时候会消失（按照考虑的精度）.这样一来，可以得到
\begin{equation*}
  \left(R^{ik}-\frac{1}{2}g^{ik}R\right)^{(2)}
  =-\frac{8\pi k}{c^{4}}\langle t^{ik(2)}\rangle
\end{equation*}
或者，由于 $\langle t_i^{i(2)}\rangle=0$，按照同样的精度：
\begin{equation*}
  \langle R_{ik}^{(2)}\rangle
  =-\frac{8\pi k}{c^{4}}\langle t_{ik}^{(2)}\rangle.
\end{equation*}

最后，采用 (108.9)，我们最终得到方程 (108.10) 的下列形式：
\begin{equation}
  R_{ik}^{(0)}
  =\frac{1}{4}\langle h_{q,i}^{n}h_{n,k}^{q}\rangle.
\end{equation}

如果“背景”完全由波本身建立，那么方程 (108.7) 和 (108.11) 应该联立求解.对方程 (108.11) 左右两边表达式的估算显示，在这种情况下背景度规的曲率半径的数量级 $L$，波长 $\lambda$ 以及它的场的数量级 $h$ 之间可以按照关系式 $L^{-2}\sim h^{2}/\lambda^{2}$，即 $\lambda/L\sim h$ 联系起来.

\section{强引力波}

本节将研究爱因斯坦方程的一种解，这种解是平直时空里的弱平面引力波的推广 (I. Robinson, H. Bondi, 1957).

我们将寻找这样的解，在其中度规张量的所有分量在适当的参考系选取下仅仅表现为单个变量的函数，我们把这个变量称做 $x^{0}$（但是并不预先确定它的特点）.这个条件允许进行下列形式的坐标变换：
\begin{equation}
  x^{\alpha}\to x^{\alpha}+\varphi^{\alpha}(x^{0}),
\end{equation}

\begin{equation}
  x^{0}\to\varphi^{0}(x^{0}),
\end{equation}
其中 $\varphi^0$，$\varphi^\alpha$ 为任意的函数.

这个解的特性本质上依赖于我们能否通过三个变换 (109.1) 使得所有的 $g_{0\alpha}$ 归零.可以做到这一点的条件是行列式 $|g_{\alpha\beta}|\neq0$.实际上，在进行变换 (109.1) 时 $g_{0\alpha}\to g_{0\alpha}+g_{\alpha\beta}\dot{\varphi}^{\beta}$（其中符号上方的点代表对 $x^{0}$ 微分）；当 $|g_{\alpha\beta}|\neq0$ 时方程组
\begin{equation*}
  g_{0\alpha}+g_{\alpha\beta}\dot{\varphi}^{\beta}=0
\end{equation*}
就确定了能够实现所要求变换的函数 $\varphi^{\beta}(x^{0})$.这种情况将会在§117中进行研究；这里我们只对这样的解感兴趣，其中
\begin{equation}
  |g_{\alpha\beta}|=0.
\end{equation}

在这种情况下不存在这样的参考系，在其中所有的 $g_{0\alpha}=0$.然而，作为替代，采用 4 个变换 (109.1)，(109.2) 可以使得下列各式得到满足：
\begin{equation}
  g_{01}=1,\qquad g_{00}=g_{02}=g_{03}=0.
\end{equation}

在这种条件下变量 $x^{0}$ 具有“类光”的特征：当 $\mathrm{d}x^{\alpha}=0$，$\mathrm{d}x^{0}\neq0$ 时，间隔 $\mathrm{d}s=0$；以这种方式选取的变量 $x^{0}$ 在下文中将表示为 $x^{0}=\eta$.在条件 (109.4) 下，线元可以表示为如下形式：
\begin{equation}
  \mathrm{d}s^{2}=2\,\mathrm{d}x^{1}\mathrm{d}\eta
  +g_{ab}(\mathrm{d}x^{a}+g^{a}\mathrm{d}x^{1})
  (\mathrm{d}x^{b}+g^{b}\mathrm{d}x^{1}).
\end{equation}

在本节里，此处和下文中指标 $a,b,c,\cdots$ 的取值范围为 2, 3；$g_{ab}(\eta)$ 可以看做二维张量，而 $g^{a}(\eta)$ 这两个量是二维矢量的分量.量 $R_{ab}$ 的计算会导致下面的场方程
\begin{equation*}
  R_{ab}=-\frac{1}{2}g_{ac}\dot{g}^{c}g_{bd}\dot{g}^{d}=0.
\end{equation*}

由此可以得出，$g_{ac}\dot{g}^{c}=0$ 或者 $\dot{g}^{c}=0$，即 $g^{c}=\mathrm{const}$.利用变换 $x^{a}+g^{a}x^{1}\to x^{a}$ 可以将被研究的度规转化成下面的形式：
\begin{equation}
  \mathrm{d}s^{2}=2\,\mathrm{d}x^{1}\mathrm{d}\eta
  +g_{ab}(\eta)\,\mathrm{d}x^{a}\mathrm{d}x^{b}.
\end{equation}

这个度规张量的行列式 $-g$ 与行列式 $|g_{ab}|$ 相同.在所有的克里斯托夫符号中只有下面的几个不为零：
\begin{equation*}
  \Gamma_{b0}^{a}=\frac{1}{2}\varkappa_b^a,
  \qquad
  \Gamma_{ab}^{1}=-\frac{1}{2}\varkappa_{ab},
\end{equation*}
其中我们引入了二维张量 $\varkappa_{ab}=\dot{g}_{ab}$，$\varkappa_a^b=g^{bc}\varkappa_{ac}$.从里奇张量的所有分量中只有 $R_{00}$ 不恒等于零，于是我们有方程
\begin{equation}
  R_{00}=-\frac{1}{2}\dot{\varkappa}_a^a
  -\frac{1}{4}\varkappa_a^b\varkappa_b^a=0.
\end{equation}

这样一来，三个函数 $g_{22}(\eta)$，$g_{23}(\eta)$，$g_{33}(\eta)$ 总共只应该满足一个方程.因此它们其中的两个可以任意给定.为方便起见，将方程 (109.7) 表示成另一种形式.首先将 $g_{ab}$ 写成下面的形式：
\begin{equation}
  g_{ab}=-\chi^{2}\gamma_{ab},\qquad|\gamma_{ab}|=1.
\end{equation}

于是行列式 $-g=|g_{ab}|=\chi^{4}$，将其代入 (109.7)，在简单的变换之后给出
\begin{equation}
  \ddot{\chi}
  +\frac{1}{8}(\dot{\gamma}_{ac}\gamma^{bc})
  (\dot{\gamma}_{bd}\gamma^{ad})\chi=0
\end{equation}
（$\gamma^{ab}$ 是二维张量，是 $\gamma_{ab}$ 的逆张量）.如果给定任意的函数 $\gamma_{ab}(\eta)$（相互之间通过关系式 $|\gamma_{ab}|=1$ 联系），就能由这些方程确定函数 $\chi(\eta)$.

这样一来我们得到包含两个任意函数的解.容易看出，这个解是§107研究的沿着一个方向传播的弱平面引力波的推广\footnote{变量数更多、性质类似的解可以参考 I. Robinson, A. Trautman//Phys. Rev. Lett. 1960. V. 4. P. 431; Proc. Roy. Soc. 1962. V. A265. P. 463.}.如果进行下面的变换
\begin{equation*}
  \eta=\frac{t+x}{\sqrt{2}},
  \qquad
  x^{1}=\frac{t-x}{\sqrt{2}},
\end{equation*}
并且令 $\gamma_{ab}=\delta_{ab}+h_{ab}(\eta)$（其中 $h_{ab}$ 是符合条件 $h_{22}+h_{33}=0$ 的小量）和 $\chi=1$，就可以得到弱平面引力波；如果忽略其中的二阶小量的项，则常数值 $\chi$ 满足方程 (109.9).

假设有限尺度的弱引力波（“波包”）通过空间的某一点.在开始通过之前我们有 $h_{ab}=0$，$\chi=1$；在引力波完全通过之后再次有 $h_{ab}=0$，$\partial^{2}\chi/\partial t^{2}=0$.但是考虑方程 (109.9) 中的二次项会导致出现不为零的负值 $\partial\chi/\partial t$：
\begin{equation*}
  \frac{\partial\chi}{\partial t}\approx
  -\frac{1}{8}\int\left(\frac{\partial h_{ab}}{\partial t}\right)^{2}
  \mathrm{d}t<0
\end{equation*}
（积分范围是波的通过时间）.因此在波通过之后将有 $\chi=1-\mathrm{const}\cdot t$，并且经过一有限时间间隔后 $\chi$ 会变号.但 $\chi$ 变为零就是度规行列式 $g$ 变为零，即度规中的奇异性.然而这个奇异性在本质上并不是物理的；它只与被通过的引力波“损坏”的参考系的缺陷有关，并且这个缺陷可以通过适当的参考系变换来修复；在引力波通过之后时空实际上又重新成为平直的.

这一点可以直接予以证明.如果变量 $\eta$ 从它对应于奇点的值起测度，那么 $\chi=\eta$，于是
\begin{equation*}
  \mathrm{d}s^{2}=2\,\mathrm{d}\eta\,\mathrm{d}x^{1}
  -\eta^{2}\left[(\mathrm{d}x^{2})^{2}+(\mathrm{d}x^{3})^{2}\right].
\end{equation*}

作变换
\begin{equation*}
  \eta x^{2}=y,\qquad
  \eta x^{3}=z,\qquad
  x^{1}=\xi-\frac{y^{2}+z^{2}}{2\eta}
\end{equation*}
之后我们得到
\begin{equation*}
  \mathrm{d}s^{2}=2\,\mathrm{d}\eta\,\mathrm{d}\xi
  -\mathrm{d}y^{2}-\mathrm{d}z^{2},
\end{equation*}
再作代换 $\eta=(t+x)/\sqrt{2}$，$\xi=(t-x)/\sqrt{2}$，最终导出伽利略形式的度规.

引力波的这个性质（虚假奇异性的产生）当然与波的微弱性无关，它是方程 (109.7) 的通解本质上所具有的性质；就像已经研究的例子，在奇异性附近 $\chi\sim\eta$，即 $-g\sim\eta^{4}$.\footnote{可以完全借鉴 §97 节中在同步参考系内类似的三维方程的方法，借助于方程（109.7）来说明这一点. 如同那里所说，虚假奇异性的产生与坐标线的交叉有关.}

\begin{xiti}

\noindent{\heiti 习题}\quad 求使下面形式的度规
\begin{equation*}
  \mathrm{d}s^{2}=\mathrm{d}t^{2}-\mathrm{d}x^{2}-\mathrm{d}y^{2}-\mathrm{d}z^{2}
  +f(t-x,y,z)(\mathrm{d}t-\mathrm{d}x)^{2}
\end{equation*}
成为真空中的爱因斯坦场方程严格解的条件 (A. Peres, 1960).

解：在坐标 $u=(t-x)/\sqrt{2}$，$v=(t+x)/\sqrt{2}$，$y,z$ 中最容易计算里奇张量，在其中
\begin{equation*}
  \mathrm{d}s^{2}=-\mathrm{d}y^{2}-\mathrm{d}z^{2}
  +2\,\mathrm{d}u\,\mathrm{d}v+2f(u,y,z)\,\mathrm{d}u^{2}.
\end{equation*}

除了 $g_{22}=g_{33}=-1$ 之外，只有下列度规张量的分量不为零：$g_{uu}=2f$，$g_{uv}=1$；在这种情况下 $g^{vv}=-2f$，$g^{uv}=1$，而行列式 $g=-1$.按照 (92.1) 进行直接计算，针对不为零的曲率张量的分量，给出下列结果：
\begin{equation*}
  R_{yuyu}=-\frac{\partial^{2}f}{\partial y^{2}},
  \qquad
  R_{zuzu}=-\frac{\partial^{2}f}{\partial z^{2}},
  \qquad
  R_{yuzu}=-\frac{\partial^{2}f}{\partial y\,\partial z}.
\end{equation*}

里奇张量唯一的非零分量是：$R_{uu}=\Delta f$，其中 $\Delta$ 是关于坐标 $y,z$ 的拉普拉斯算符.这样一来，爱因斯坦方程是：$\Delta f=0$，就是说，函数 $f(t-x,y,z)$ 应该是关于变量 $y,z$ 的调和函数.

如果函数 $f$ 不依赖于 $y,z$，或者和这些变量是线性关系，那么场就不存在——时空是平直的（曲率张量为零）.关于 $y,z$ 的二次函数
\begin{equation*}
  f(u,y,z)=yzf_1(u)+\frac{1}{2}(y^{2}-z^{2})f_2(u)
\end{equation*}
对应向 $x$ 轴的正方向传播的平面波；实际上，曲率张量在这样的场中只依赖于 $t-x$：
\begin{equation*}
  R_{yuzu}=-f_1(u),
  \qquad
  R_{yuyu}=-R_{zuzu}=-f_2(u).
\end{equation*}

对应于波的两个可能的偏振，在这种情况下度规包含两个任意的函数 $f_1(u)$ 和 $f_2(u)$.

\end{xiti}
\section{引力波的辐射}

我们下面来研究一个运动速度比光速小很多的物体所产生的弱引力场.

由于物质的存在，引力场方程将不同于简单的波动方程 $\square h_i^k=0$ (107.8)，其差异在于等式右边有来自物质的能量动量张量的项.我们将这些方程写成
\begin{equation}
  \frac{1}{2}\square\psi_i^k=\frac{8\pi k}{c^{4}}\tau_i^k,
\end{equation}
其中我们引入了对于这种情形更便利的量
\begin{equation*}
  \psi_i^k=h_i^k-\frac{1}{2}\delta_i^kh
\end{equation*}
来代替 $h_i^k$，而 $\tau_i^k$ 则用来标记辅助量，从严格的引力方程出发，作弱场近似就会得到这些量.不难证明，分量 $\tau_0^0$ 和 $\tau_\alpha^0$ 可以直接从相应的分量 $T_i^k$ 得来，只需从 $T_i^k$ 中取出我们感兴趣的量级的量即可；至于分量 $\tau_\beta^\alpha$，它们除了包含从 $T_\beta^\alpha$ 得来的项外，还包含从 $R_i^k-\delta_i^kR/2$ 得来的二级小量的项\footnote{从方程（110.1）可以再次得到在 §106 节中使用过的针对物体远处的弱恒定场公式（106.1）和（106.2）. 在一级近似中，我们可以略去含有对时间的二阶导数的那些项（含有 $1/c^{2}$ 的那些项），而在 $\tau_{i}^{k}$ 的所有分量之中只保留 $\tau_{0}^{0}=\mu c^{2}$. 方程 $\Delta\psi_{\alpha}^{\beta}=0$，$\Delta\psi_{0}^{\alpha}=0$，$\Delta\psi_{0}^{0}=16\pi k\mu/c^{2}$ 在无穷远处变为零的解为 $\psi_{\alpha}^{\beta}=0$，$\psi_{0}^{\alpha}=0$，$\psi_{0}^{0}=4\varphi/c^{2}$，其中 $\varphi$ 是牛顿引力势，参考方程（99.2）. 由此对于张量 $h_{i}^{k}=\psi_{i}^{k}-\psi\delta_{i}^{k}/2$ 得到值（106.1），（106.2）.}.

$\psi_i^k$ 满足条件 (107.5) $\partial\psi_i^k/\partial x^k=0$.从 (110.1) 可以推断，同样的方程对于 $\tau_i^k$ 也成立：
\begin{equation}
  \frac{\partial\tau_i^k}{\partial x^{k}}=0.
\end{equation}

这个方程在这里就代替了普遍关系式 $T_{i;k}^k=0$.

借助于得到的方程，我们来研究运动的物体以引力波的形式辐射的能量.这个问题的解决需要确定在“波区”的引力场，就是在距离远大于辐射波波长之处的引力场.

原则上，所有的计算完全与对电磁波所作的计算相似.弱引力场的方程 (110.1) 在形式上与推迟势的方程（§62）完全一样.因此，我们立刻能够写出它的通解如下：
\begin{equation}
  \psi_i^k=-\frac{4k}{c^{4}}\int(\tau_i^k)_{t-R/c}
  \frac{\mathrm{d}V}{R}.
\end{equation}

既然体系内的所有物体的速度很小，那么，我们就能够写出与体系相距甚远之处的场（见§66 和§67）：
\begin{equation}
  \psi_i^k=-\frac{4k}{c^{4}R_0}\int(\tau_i^k)_{t-R_0/c}\,\mathrm{d}V,
\end{equation}
其中，$R_0$ 是到原点的距离，原点被选择在体系内的任意一点.为简单起见，此后我们将略去被积函数内的脚标 $t-R_0/c$.

为了计算这些积分，我们利用方程 (110.2).降低 $\tau_i^k$ 的指标，分开空间和时间分量，我们将 (110.2) 写成
\begin{equation}
  \frac{\partial\tau_{\alpha\gamma}}{\partial x^{\gamma}}
  -\frac{\partial\tau_{\alpha0}}{\partial x^{0}}=0,
  \qquad
  \frac{\partial\tau_{0\gamma}}{\partial x^{\gamma}}
  -\frac{\partial\tau_{00}}{\partial x^{0}}=0.
\end{equation}

用 $x^{\beta}$ 乘第一式，然后对整个空间积分，则得
\begin{equation*}
  \frac{\partial}{\partial x^{0}}\int\tau_{\alpha0}x^{\beta}\,\mathrm{d}V
  =\int\frac{\partial\tau_{\alpha\gamma}}{\partial x^{\gamma}}
  x^{\beta}\,\mathrm{d}V
  =\int\frac{\partial(\tau_{\alpha\gamma}x^{\beta})}{\partial x^{\gamma}}
  \,\mathrm{d}V-\int\tau_{\alpha\beta}\,\mathrm{d}V.
\end{equation*}

因为在无穷远处 $\tau_{ik}=0$，右边的第一个积分在经过高斯定理的变换后消失.将余下来的方程与它交换指标后得到的方程相加，再取半，我们便求得
\begin{equation*}
  \int\tau_{\alpha\beta}\,\mathrm{d}V
  =-\frac{1}{2}\frac{\partial}{\partial x^{0}}
  \int(\tau_{\alpha0}x^{\beta}+\tau_{\beta0}x^{\alpha})\,\mathrm{d}V.
\end{equation*}

接下来，用 $x^{\alpha}x^{\beta}$ 乘 (110.5) 中的第二个方程，然后再对整个空间积分.相似的变换导出
\begin{equation*}
  \frac{\partial}{\partial x^{0}}\int\tau_{00}x^{\alpha}x^{\beta}\,\mathrm{d}V
  =-\int(\tau_{\alpha0}x^{\beta}+\tau_{\beta0}x^{\alpha})\,\mathrm{d}V.
\end{equation*}

将所得的两个结果加以比较，我们求得
\begin{equation}
  \int\tau_{\alpha\beta}\,\mathrm{d}V
  =\frac{1}{2}\left(\frac{\partial}{\partial x^{0}}\right)^{2}
  \int\tau_{00}x^{\alpha}x^{\beta}\,\mathrm{d}V.
\end{equation}

因此，所有 $\tau_{\alpha\beta}$ 的积分可以表为仅包含分量 $\tau_{00}$ 的积分.但是这些分量就像上面指出的那样，等于能量动量张量的相应的分量 $T_{00}$，并且可以以足够的精度写出（参考 (99.1)）：
\begin{equation}
  \tau_{00}=\mu c^{2}.
\end{equation}

将它代入 (110.6)，并且引入时间 $t=x^{0}/c$，将 (110.4) 改写成如下形式：
\begin{equation}
  \psi_{\alpha\beta}=-\frac{2k}{c^{4}R_0}
  \frac{\partial^{2}}{\partial t^{2}}
  \int\mu x^{\alpha}x^{\beta}\,\mathrm{d}V.
\end{equation}

在与物体体系相距甚远之处，我们可以认为波（在不大的空间区域内）是平面波.因此，利用 (107.12) 式我们可以计算出体系辐射的能流，例如沿着 $x^{1}$ 轴方向的能流.在这个公式内只包含分量 $h_{23}=\psi_{23}$ 和 $h_{22}-h_{33}=\psi_{22}-\psi_{33}$.从 (110.8)，我们求出它们的表达式：\footnote{张量（110.8）不满足推导出公式（107.12）的那些条件. 然而将 $h_{ik}$ 化为所要求的规范形式的参考系变换不会影响到这里使用的分量（110.9）的值.}
\begin{equation}
  h_{23}=-\frac{2k}{3c^{4}R_0}\ddot{D}_{23},
  \qquad
  h_{22}-h_{33}
  =-\frac{2k}{3c^{4}R_0}(\ddot{D}_{22}-\ddot{D}_{33})
\end{equation}
（符号上的一点表示对时间微分），此处我们引入了质量的四极矩张量 (99.8)：
\begin{equation}
  D_{\alpha\beta}
  =\int\mu(3x^{\alpha}x^{\beta}-r^{2}\delta_{\alpha\beta})\,\mathrm{d}V.
\end{equation}

结果我们求得沿着 $x^{1}$ 轴的能流如下：
\begin{equation}
  ct^{10}=\frac{k}{36\pi c^{5}R_0^{2}}
  \left[\left(\frac{\dddot{D}_{22}-\dddot{D}_{33}}{2}\right)^{2}
  +\dddot{D}_{23}^{2}\right].
\end{equation}

在该方向单位立体角上的能流可通过对上式乘以 $R_0^{2}\,\mathrm{d}o$ 获得.

这个表达式中的两项对应于两个独立偏振的波的辐射.为了将它们写成不变的形式（不依赖于辐射方向的选择），我们引入平面引力波的三维单位极化张量 $e_{\alpha\beta}$，这个张量确定分量 $h_{\alpha\beta}$ 中到底哪些不为零（在 $h_{ik}$ 的这个规范中，$h_{0\alpha}=h_{00}=h=0$）.极化张量是对称的，并且满足条件
\begin{equation}
  e_{\alpha\alpha}=0,\qquad
  e_{\alpha\beta}n_{\beta}=0,\qquad
  e_{\alpha\beta}e_{\alpha\beta}=1,
\end{equation}
其中 $\boldsymbol{n}$ 是沿波的传播方向的单位矢量；头两个条件表达了波的张量性和横波特性.

借助于这个张量，立体角 $\mathrm{d}o$ 上给定偏振的辐射强度可写成如下形式：
\begin{equation}
  \mathrm{d}I=\frac{k}{72\pi c^{5}}
  (\dddot{D}_{\alpha\beta}e_{\alpha\beta})^{2}\,\mathrm{d}o.
\end{equation}

由于横波条件 $e_{\alpha\beta}n_{\beta}=0$，这个表达式隐含地依赖于方向 $\boldsymbol{n}$.所有偏振的总的角分布可以通过对 (110.13) 按偏振求和来获得，或者等价于对偏振进行平均，再将结果乘以 2（独立偏振的数目）.求平均通过下面的公式实现：
\begin{equation}
\begin{aligned}
  \overline{e_{\alpha\beta}e_{\gamma\delta}}
  =\frac{1}{4}\bigl\{&n_{\alpha}n_{\beta}n_{\gamma}n_{\delta}
  +(n_{\alpha}n_{\beta}\delta_{\gamma\delta}
  +n_{\gamma}n_{\delta}\delta_{\alpha\beta})\\
  &-(n_{\alpha}n_{\gamma}\delta_{\beta\delta}
  +n_{\beta}n_{\gamma}\delta_{\alpha\delta}
  +n_{\alpha}n_{\delta}\delta_{\beta\gamma}
  +n_{\beta}n_{\delta}\delta_{\alpha\gamma})\\
  &-\delta_{\alpha\beta}\delta_{\gamma\delta}
  +(\delta_{\alpha\gamma}\delta_{\beta\delta}
  +\delta_{\beta\gamma}\delta_{\alpha\delta})\bigr\}
\end{aligned}
\end{equation}
（式子的右边是由单位张量和矢量 $\boldsymbol{n}$ 的分量组成的张量，它具备所要求的指标对称性，按指标对 $\alpha,\gamma$ 和 $\beta,\delta$ 进行缩并时给出 1，在与 $\boldsymbol{n}$ 求标量积后变为零）.结果我们得到：
\begin{equation}
  \mathrm{d}I=\frac{k}{36\pi c^{5}}
  \left[\frac{1}{4}(\dddot{D}_{\alpha\beta}n_{\alpha}n_{\beta})^{2}
  +\frac{1}{2}\dddot{D}_{\alpha\beta}^{2}
  -\dddot{D}_{\alpha\beta}\dddot{D}_{\alpha\gamma}n_{\beta}n_{\gamma}\right]
  \mathrm{d}o.
\end{equation}

沿所有方向的总辐射，即系统在单位时间内的能量损失 $(-\mathrm{d}\mathcal{E}/\mathrm{d}t)$，可以通过将 $\mathrm{d}I/\mathrm{d}o$ 对所有方向 $\boldsymbol{n}$ 求平均值，然后将所得的结果乘以 $4\pi$ 得出.利用§71第 1 个脚注中的公式，就很容易进行求平均值的计算.得到能量损失的公式如下（A. 爱因斯坦，1918）：
\begin{equation}
  -\frac{\mathrm{d}\mathcal{E}}{\mathrm{d}t}
  =\frac{k}{45c^{5}}\dddot{D}_{\alpha\beta}^{2}.
\end{equation}

我们指出，引力波的辐射是关于 $1/c$ 的五次方效应.一般说来，这一事实与微小的引力常数 $k$ 一起导致这种效应是极其微弱的.

\begin{xiti}

\noindent{\heiti 习题1}\quad 两个物体按照牛顿定律相互吸引，并绕着共同的惯性中心作圆周运动.求引力波辐射的平均强度（在一个转动周期内）及其偏振和方向的分布.

解：选取惯性中心作为坐标原点，对于两个物体的径矢，我们有：
\begin{equation*}
  \boldsymbol{r}_1=\frac{m_2}{m_1+m_2}\boldsymbol{r},
  \qquad
  \boldsymbol{r}_2=-\frac{m_1}{m_1+m_2}\boldsymbol{r},
  \qquad
  \boldsymbol{r}=\boldsymbol{r}_1-\boldsymbol{r}_2.
\end{equation*}

张量 $D_{\alpha\beta}$ 的分量是（假设 $xy$ 与运动平面重合）：
\begin{equation*}
\begin{aligned}
  &D_{xx}=\mu r^{2}(3\cos^{2}\psi-1),
  \qquad
  D_{yy}=\mu r^{2}(3\sin^{2}\psi-1),\\
  &D_{xy}=3\mu r^{2}\cos\psi\sin\psi,
  \qquad
  D_{zz}=-\mu r^{2},
\end{aligned}
\end{equation*}
其中 $\mu=m_1m_2/(m_1+m_2)$，$\psi$ 是矢量 $\boldsymbol{r}$ 在 $xy$ 平面上的极角.在圆周运动时 $r=\mathrm{const}$，
\begin{equation*}
  \dot{\psi}=r^{-3/2}\sqrt{k(m_1+m_2)}\equiv\omega.
\end{equation*}

利用球面角（极角 $\theta$ 和方位角 $\varphi$）和垂直于运动平面的极轴 $z$ 来给定方向 $\boldsymbol{n}$.我们考察两个偏振，对于它们：1) $e_{\theta\varphi}=1/\sqrt{2}$；2) $e_{\theta\theta}=-e_{\varphi\varphi}=1/\sqrt{2}$.

将张量 $D_{\alpha\beta}$ 投影到球面单位矢量 $e_\theta$ 和 $e_\varphi$ 的方向上，按照公式 (110.13) 计算并对时间求平均，结果对这两种情况以及总合 $I=I_1+I_2$ 我们得到：
\begin{equation*}
  \frac{\overline{\mathrm{d}I_1}}{\mathrm{d}o}
  =\frac{k\mu^{2}\omega^{6}r^{4}}{2\pi c^{5}}\cdot4\cos^{2}\theta,
  \qquad
  \frac{\overline{\mathrm{d}I_2}}{\mathrm{d}o}
  =\frac{k\mu^{2}\omega^{6}r^{4}}{2\pi c^{5}}(1+\cos^{2}\theta)^{2},
\end{equation*}

\begin{equation*}
  \frac{\overline{\mathrm{d}I}}{\mathrm{d}o}
  =\frac{k\mu^{2}\omega^{6}r^{4}}{2\pi c^{5}}
  (1+6\cos^{2}\theta+\cos^{4}\theta),
\end{equation*}
然后沿方向积分后：
\begin{equation*}
  -\frac{\mathrm{d}\mathcal{E}}{\mathrm{d}t}
  =I=\frac{32k\mu^{2}\omega^{6}r^{4}}{5c^{5}}
  =\frac{32k^{4}m_1^{2}m_2^{2}(m_1+m_2)}{5c^{5}r^{5}},
  \qquad
  \frac{\overline{I}_1}{\overline{I}_2}=\frac{5}{7}
\end{equation*}
（若只计算总的强度 $I$，当然应当使用 (110.16)）.

辐射系统的能量损失导致两个物体逐渐（长期的）靠近.因为 $\mathcal{E}=-km_1m_2/2r$，那么靠近速度是
\begin{equation*}
  \dot{r}=\frac{2r^{2}}{km_1m_2}\frac{\mathrm{d}\mathcal{E}}{\mathrm{d}t}
  =-\frac{64k^{3}m_1m_2(m_1+m_2)}{5c^{5}r^{3}}.
\end{equation*}

\noindent{\heiti 习题2}\quad 求两个沿椭圆轨道运动的物体组成的系统以引力波形式辐射的平均能量（对一个转动周期求平均）(P. C. Peters, J. Mathews).\footnote{关于这个辐射的角分布，偏振分布和谱分布可参考 Phys. Rev. 1963. V. 131. P. 435.}

解：区别于圆周运动情况，距离 $r$ 和角速度沿着轨道按下面的规律变化：
\begin{equation*}
  \frac{a(1-e^{2})}{r}=1+e\cos\psi,
  \qquad
  \frac{\mathrm{d}\psi}{\mathrm{d}t}
  =\frac{1}{r^{2}}\left[k(m_1+m_2)a(1-e^{2})\right]^{1/2},
\end{equation*}
其中 $e$ 是偏心率，而 $a$ 是轨道的半长轴（参考本教程第一卷§15）.使用 (110.16) 进行相当长的运算得到
\begin{equation*}
  -\frac{\mathrm{d}\mathcal{E}}{\mathrm{d}t}
  =\frac{8k^{4}m_1^{2}m_2^{2}(m_1+m_2)}
  {15a^{5}c^{5}(1-e^{2})^{5}}
  (1+e\cos\psi)^{4}
  \left[12(1+e\cos\psi)^{2}+e^{2}\sin^{2}\psi\right].
\end{equation*}

对一个转动周期求平均时使用对 $\mathrm{d}\psi$ 的积分替换对 $\mathrm{d}t$ 的积分，导出如下结果：
\begin{equation*}
  -\frac{\overline{\mathrm{d}\mathcal{E}}}{\mathrm{d}t}
  =\frac{32k^{4}m_1^{2}m_2^{2}(m_1+m_2)}{5c^{5}a^{5}}
  \frac{1}{(1-e^{2})^{7/2}}
  \left(1+\frac{73}{24}e^{2}+\frac{37}{96}e^{4}\right).
\end{equation*}

我们注意到辐射强度随着轨道偏心率的增加而快速增长.

\noindent{\heiti 习题3}\quad 由稳态运动的物体组成的系统辐射引力波，求该系统角动量的平均（按照时间）损失速率.

解：为了便于书写公式，我们将物体系统临时看做是由一些离散的粒子组成的.系统能量的平均损失速率可以认为是作用在粒子上的“摩擦力”$\boldsymbol{f}$ 所做的功：
\begin{equation}
  \frac{\overline{\mathrm{d}\mathcal{E}}}{\mathrm{d}t}
  =\sum\overline{\boldsymbol{f}\cdot\boldsymbol{v}}
  \tag{1}
\end{equation}
（给粒子编号的指标没有写出）.那么角动量的平均损失速率可以这样计算：
\begin{equation}
  \overline{\frac{\mathrm{d}M_{\alpha}}{\mathrm{d}t}}
  =\sum\overline{(\boldsymbol{r}\times\boldsymbol{f})_{\alpha}}
  =\sum e_{\alpha\beta\gamma}\overline{x_{\beta}f_{\gamma}}
  \tag{2}
\end{equation}
（与公式 (75.7) 的推导进行比较）.为了确定 $\boldsymbol{f}$ 我们写出
\begin{equation*}
  \frac{\overline{\mathrm{d}\mathcal{E}}}{\mathrm{d}t}
  =-\frac{k}{45c^{5}}\overline{\dddot{D}_{\alpha\beta}
  \dddot{D}_{\alpha\beta}}
  =-\frac{k}{45c^{5}}\overline{\dot{D}_{\alpha\beta}
  \overset{\cdots\cdots}{D}_{\alpha\beta}}
\end{equation*}
（使用了对时间的全导数的平均值等于零这个等式）.将
$\dot{D}_{\alpha\beta}=\sum m(3x_{\alpha}v_{\beta}
+3x_{\beta}v_{\alpha}-2\boldsymbol{r}\cdot\boldsymbol{v}\,\delta_{\alpha\beta})$
代入此式并且和 (1) 进行比较，我们得到
\begin{equation*}
  f_{\alpha}=-\frac{2k}{15c^{5}}\ddot{D}_{\alpha\beta}\,mx_{\beta}.
\end{equation*}

将上式代入 (2) 可导出结果：
\begin{equation}
  \overline{\frac{\mathrm{d}M_{\alpha}}{\mathrm{d}t}}
  =-\frac{2k}{45c^{5}}e_{\alpha\beta\gamma}
  \overline{\dddot{D}_{\beta\delta}D_{\gamma\delta}}
  \tag{3}
\end{equation}

\noindent{\heiti 习题4}\quad 针对两个沿椭圆轨道运动的物体组成的系统，求其在单位时间内平均损失的角动量.

解：使用上一习题的公式 (3) 进行计算，类似于习题2 中的推导，可得到结果：
\begin{equation*}
  -\frac{\overline{\mathrm{d}M_z}}{\mathrm{d}t}
  =\frac{32k^{7/2}m_1^{2}m_2^{2}\sqrt{m_1+m_2}}{5c^{5}a^{7/2}}
  \frac{1}{(1-e^{2})^{2}}
  \left(1+\frac{7}{8}e^{2}\right).
\end{equation*}

对于圆周运动（$e=0$），$\dot{\mathcal{E}}$ 和 $\dot{M}$ 的值由 $\dot{\mathcal{E}}=\dot{M}\omega$ 联系，这也是它们之间应有的关系式.

\end{xiti}
